\begin{proof}[Proof of \cref{obj:thm:fixed-binary-minimax}]
\begin{enumerate}
\item Define the two testing experiments, their target values, and their observed laws by
\[
\begin{aligned}
\mathsf M_\pm
&:=\text{the experiment of \cref{obj:def:four-state-testing-family}
with }v=\pm v_T,\\
\theta_\pm&:=\theta(K_{\pm v_T}^{(T)},e),\\
\mathbb P_\pm^{\mathsf O_T}
&:=\mathcal L_{\mathsf M_\pm}(\mathsf O_T).
\end{aligned}
\]
By \cref{obj:lem:testing-family-membership},
\[
\mathsf M_\pm\in\mathcal M_T^{(2)}(t_0,\zeta),
\qquad
|\theta_+-\theta_-|=\frac1{8\sqrt T},
\qquad
\operatorname{KL}\!\left(
\mathbb P_+^{\mathsf O_T}\Vert\mathbb P_-^{\mathsf O_T}
\right)\leq\frac1{12},
\]
and the relative entropy is finite. Both experiments supply the same fair behavior and target policies to the estimator.
% lean: CausalSmith.Stat.PomdpBinaryhiddenRate.testingFamily_membership

\item Fix an admissible estimator \(\widetilde\theta\) from
\cref{obj:def:minimax-risk}. Because its supplied policies agree across the pair, it is the same measurable function of the observed trajectory in both experiments. Pinsker's inequality gives
\[
\left\|\mathbb P_+^{\mathsf O_T}
-\mathbb P_-^{\mathsf O_T}\right\|_{\mathrm{TV}}
\leq
\sqrt{\frac{
\operatorname{KL}(\mathbb P_+^{\mathsf O_T}
\Vert\mathbb P_-^{\mathsf O_T})}{2}}
\leq\sqrt{\frac1{24}}\leq\frac14.
\]
Define the error threshold and the two measurable error events by
\[
\varepsilon_{\mathrm{test}}:=\frac1{16\sqrt T},
\qquad
\mathcal E_\pm:=
\left\{|\widetilde\theta-\theta_\pm|
\geq\varepsilon_{\mathrm{test}}\right\}.
\]
Then
\[
|\theta_+-\theta_-|=2\varepsilon_{\mathrm{test}},
\qquad
\mathcal E_+^{\mathsf c}\subseteq\mathcal E_-,
\]
where the inclusion follows from the triangle inequality. Consequently,
\[
\begin{aligned}
\mathbb P_+^{\mathsf O_T}(\mathcal E_+)
+\mathbb P_-^{\mathsf O_T}(\mathcal E_-)
&\geq
\mathbb P_+^{\mathsf O_T}(\mathcal E_+)
+\mathbb P_-^{\mathsf O_T}(\mathcal E_+^{\mathsf c})\\
&\geq
1-\left\|\mathbb P_+^{\mathsf O_T}
-\mathbb P_-^{\mathsf O_T}\right\|_{\mathrm{TV}}
\geq\frac34.
\end{aligned}
\]
Thus
\[
\max\left\{
\mathbb P_+^{\mathsf O_T}(\mathcal E_+),
\mathbb P_-^{\mathsf O_T}(\mathcal E_-)
\right\}\geq\frac38.
\]
For each sign, integrating the pointwise inequality
\[
(\widetilde\theta-\theta_\pm)^2
\geq\varepsilon_{\mathrm{test}}^2\mathbf1_{\mathcal E_\pm}
\]
gives
\[
\mathbb E_{\mathsf M_\pm}
[(\widetilde\theta-\theta_\pm)^2]
\geq
\varepsilon_{\mathrm{test}}^2
\mathbb P_\pm^{\mathsf O_T}(\mathcal E_\pm).
\]
If the positive-sign error probability is at most the negative-sign error probability, apply this inequality to the negative sign; in the opposite case, apply it to the positive sign. In either case,
\[
\begin{aligned}
\max\left\{
\mathbb E_{\mathsf M_+}[(\widetilde\theta-\theta_+)^2],
\mathbb E_{\mathsf M_-}[(\widetilde\theta-\theta_-)^2]
\right\}
&\geq\frac38\varepsilon_{\mathrm{test}}^2\\
&=\frac3{2048T}
\geq\frac1{1024T},
\end{aligned}
\]
using \((\sqrt T)^2=T>0\).
% lean: CausalSmith.Stat.PomdpBinaryhiddenRate.testingPair_observedRisk_lower

\item For any admissible experiment \(\mathsf M\), define its risk by
\[
\mathcal R(\widetilde\theta,\mathsf M)
:=
\mathbb E_{\mathsf M}[(\widetilde\theta-\theta)^2],
\qquad
\mathsf M\in\mathcal M_T^{(2)}(t_0,\zeta).
\]
The target value in \cref{obj:def:target-value} is a stationary average of rewards in \([0,1]\). Since the estimator also takes values in \([0,1]\),
\[
0\leq\theta\leq1,
\qquad
0\leq(\widetilde\theta-\theta)^2\leq1,
\qquad
0\leq\mathcal R(\widetilde\theta,\mathsf M)\leq1.
\]
Thus every estimator's risks are bounded above over the class. The estimator family is nonempty, as witnessed by the constant estimator
\[
\widetilde\theta_0\equiv0.
\]
Membership of the testing pair and the preceding two-point bound therefore give
\[
\begin{aligned}
R_T^{(2)}(t_0,\zeta)
&=\inf_{\widetilde\theta}
\sup_{\mathsf M\in\mathcal M_T^{(2)}(t_0,\zeta)}
\mathcal R(\widetilde\theta,\mathsf M)\\
&\geq
\inf_{\widetilde\theta}
\max\left\{
\mathcal R(\widetilde\theta,\mathsf M_+),
\mathcal R(\widetilde\theta,\mathsf M_-)
\right\}
\geq\frac1{1024T}.
\end{aligned}
\]
% lean: CausalSmith.Stat.PomdpBinaryhiddenRate.observedRisk_unit_bounds
% lean: Causalean.Stat.le_minimaxValue_of_two_point

\item Use the stable-grid estimator of
\cref{obj:def:stable-grid-estimator}. Its selected reward vector and stationary probability vector satisfy
\[
0\leq r_*\leq\mathbf1,
\qquad
\pi(P_*)\geq0,
\qquad
\pi(P_*)\mathbf1=1.
\]
Hence
\[
0\leq\widehat\theta_T
=\pi(P_*)r_*
\leq\pi(P_*)\mathbf1=1.
\]
The finite candidate selection uses measurable comparison events, and the selected stationary reward is therefore measurable for each pair of supplied policies. This makes \(\widehat\theta_T\) admissible in
\cref{obj:def:minimax-risk}. By \cref{obj:thm:stable-grid-upper},
\[
\mathcal R(\widehat\theta_T,\mathsf M)
\leq\frac{B_\alpha^2(112V_{\alpha,L}+2)}{T}
\qquad
\text{for every }\mathsf M\in\mathcal M_T^{(2)}(t_0,\zeta).
\]
The class is nonempty because it contains \(\mathsf M_+\). Nonnegativity of the risks and the defining infimum and supremum now yield
\[
R_T^{(2)}(t_0,\zeta)
\leq
\sup_{\mathsf M\in\mathcal M_T^{(2)}(t_0,\zeta)}
\mathcal R(\widehat\theta_T,\mathsf M)
\leq\frac{B_\alpha^2(112V_{\alpha,L}+2)}{T}.
\]
% lean: CausalSmith.Stat.PomdpBinaryhiddenRate.stableGridEstimator_admissible
% lean: CausalSmith.Stat.PomdpBinaryhiddenRate.stableGrid_upper
% lean: Causalean.Stat.minimaxValue_le_worstCaseRisk_of_nonneg
% lean: Causalean.Stat.worstCaseRisk_le

\item The remaining assertions about the binary testing pair follow from
\cref{obj:lem:testing-family-membership}:
\[
e(a\mid x)=b(a\mid x)=\frac12,
\qquad d_e=d_b,
\]
and, for every \(C>1\),
\[
d_e(s)\leq C\,d_b(s)
\qquad\text{for every }s\in\mathcal S.
\]
Together with the estimator-wise two-point risk bound, these establish all the stated properties of the pair.
% lean: CausalSmith.Stat.PomdpBinaryhiddenRate.fixedBinary_minimax

\item Suppose now that \(\zeta>0\), and fix \(C>1\). Use the theorem's notation
\[
\beta=\frac2{2+t_0\zeta},
\qquad
q_C=\frac{C-1}{C}.
\]
The cited cardinality-uniform result supplies its signed-depth lower construction
\citep[Theorem 1 (Uniform overlap frontier) and the Signed-depth family definition]{CausalSmithLatentOverlap2026}.
Specifically, there are
\[
c_{\mathrm{Reset}}>0,\qquad
Q:\mathbb N\to\mathbb N,\qquad a_1,a_2>0
\]
such that, for every sufficiently large \(n\),
\[
a_1\log n\leq Q(n)\leq a_2\log n,
\qquad n\geq1,\qquad Q(n)\geq1.
\]
For these \(n\), define
\[
\begin{aligned}
\mathsf M_{n,\pm}^{\mathrm{Reset}}
&:=\text{the signed depth-}Q(n)\text{ reset experiments},\\
\theta_{n,\pm}^{\mathrm{Reset}}
&:=\text{their respective stationary target values}.
\end{aligned}
\]
Each experiment has exactly \(2(Q(n)+1)\) hidden states. For every admissible observable-data estimator
\(\widetilde\theta_{\mathrm{Reset}}\) in the reset decision problem, with values in \([-1,1]\), the cited construction gives
\[
\max\left\{
\mathbb E_{\mathsf M_{n,-}^{\mathrm{Reset}}}
[(\widetilde\theta_{\mathrm{Reset}}
-\theta_{n,-}^{\mathrm{Reset}})^2],
\mathbb E_{\mathsf M_{n,+}^{\mathrm{Reset}}}
[(\widetilde\theta_{\mathrm{Reset}}
-\theta_{n,+}^{\mathrm{Reset}})^2]
\right\}
\geq c_{\mathrm{Reset}}n^{-\beta}.
\]
Finally, the same eventual condition \(Q(n)\geq1\) gives
\[
2(Q(n)+1)\geq4>2,
\]
which proves the asserted cardinality inequality for the negative-sign experiment.
% lean: CausalSmith.Stat.PomdpBinaryhiddenRate.comparator_reset_depth_growth

\item Since \(t_0\zeta>0\),
\[
2+t_0\zeta>2>0,
\qquad
0<\beta=\frac2{2+t_0\zeta}<1,
\qquad
1-\beta>0.
\]
Also, \(C>1\) implies
\[
q_C>0,
\qquad q_C^{\,2(1-\beta)}>0.
\]
% lean: CausalSmith.Stat.PomdpBinaryhiddenRate.fixedBinary_minimax

\item Along the positive natural numbers, the negative-power limit gives
\[
n^{-(1-\beta)}\longrightarrow0.
\]
Multiplication by the fixed constant
\(q_C^{-2(1-\beta)}\) preserves this limit. For every \(n\geq1\), the laws of real powers give
\[
\frac{n^{-1}}{n^{-\beta}}
=n^{\beta-1}=n^{-(1-\beta)},
\]
and hence
\[
\frac{n^{-1}}
{n^{-\beta}q_C^{\,2(1-\beta)}}
=
n^{-(1-\beta)}q_C^{-2(1-\beta)}
\longrightarrow0.
\]
This proves the claimed limit and completes the proof.
% lean: CausalSmith.Stat.PomdpBinaryhiddenRate.fixedBinary_minimax
\end{enumerate}
\end{proof}
